\FloatBarrier
\section{Observations and Limitations}
\label{sec:discussion}

\subsection{Key observations}
Observations 1 and 4 concern the 2018-02-20 batch analysed in
Section~\ref{sec:eda}; observations 2 and 3 also hold for the paper's
batches (Section~\ref{sec:reproduction}).
\begin{enumerate}
  \item \textbf{All policies take the same charge time.} In this batch every
  policy reaches 80\% SOC in about 10\,min, so differences in lifetime
  reflect how the current is distributed in time, not how fast the cell is
  charged overall. (This is not true of every batch: in 2017-05-12 the
  charge time ranges from 9.0 to 13.4\,min.)
  \item \textbf{Capacity alone is a poor early predictor.} Fade stays small
  and non-linear until the knee, so early capacity says little about
  eventual lifetime. On the paper's training cells, the capacity at cycle
  100 correlates with log life at only $r=+0.47$.
  \item \textbf{Curve-shape features carry early information.} $\dq$ visibly
  separates short- and long-lived cells at cycle 100, long before EOL.
  \item \textbf{The test procedure dominates.} Whatever
  \texttt{-newstructure} changed roughly doubles cycle life. Its effect is
  larger than that of any charging policy. These cells also tend to start
  with lower internal resistance.
\end{enumerate}

\subsection{Limitations of the current analysis}
\begin{itemize}
  \item \textbf{Small sample, two groups.} The unseen batch has 39 labelled
  cells in two well-separated clusters, so a model can score well on it just
  by telling the groups apart, without predicting differences within a
  group. Results on this batch (RQ5) must therefore also be reported within
  each group.
  \item \textbf{Batch effects in the training data.} The 41 training cells
  come from two batches with different rest periods and different typical
  lives. Some features (temperature integral, charge time) partly encode the
  batch or the policy rather than the state of the cell
  (Section~\ref{sec:reproduction}).
  \item \textbf{The meaning of \texttt{-newstructure} is unknown.} The data
  does not record what was changed. Until this is clarified, these cells
  should be modelled separately or given an explicit indicator feature.
  \item \textbf{Noise.} Single-cycle glitches and the 20\,min charge-time
  cap are handled by rules B3--B5. The $\dq$ features still use only two
  cycles (10 and 100), so noise in either cycle passes straight into
  them.
\end{itemize}
